%%=============================================================================
%% Conclusie
%%=============================================================================

\chapter{\IfLanguageName{dutch}{Conclusie}{Conclusion}}%
\label{ch:conclusie}

Uit dit onderzoek blijkt dat het succesvol 
ontwikkelen van een gepersonaliseerd 
e-mailnotificatiesysteem voor de Vlaamse 
concertmarkt een architectuur vereist die 
zowel de vluchtigheid van online eventdata 
oplost als de eindgebruiker beschermt 
tegen notificatievermoeidheid.

Ter beantwoording van de specifieke deelvragen 
binnen het \textbf{probleemdomein}, is via de 
literatuurstudie vastgesteld dat metadata 
zoals starttijd, locatie en ticketprijs 
essentieel zijn voor een bruikbare en 
herbruikbare concertontologie \autocite{Graves2003}. 

De analyse van de bestaande diensten toonde 
daarnaast aan dat internationale platformen 
zoals Bandsintown of Songkick functioneel 
tekortschieten voor de lokale markt. Dit komt 
doordat zij zich voornamelijk richten op de 
mainstream en onvoldoende dekking 
bieden voor de specifieke Vlaamse zalen. 

Omdat deze lokale zalen hun kalenders zeer 
dynamisch en onvoorspelbaar aanpassen, en er 
geen uniforme API's beschikbaar zijn, vormt 
geautomatiseerde webscraping de enige effectieve 
methode om deze data tijdig te verzamelen. 

Om te voorkomen dat deze continue datastroom 
leidt tot informatie-overload, worden 
matches gebundeld en verstuurd als
één gepersonaliseerde notificatie per 24 uur.

Wat betreft de deelvragen binnen het 
\textbf{oplossingsdomein}, heeft de vergelijkende 
toolselectie aangetoond dat het \texttt{Scrapling}-framework 
het meest geschikt is voor periodieke monitoring. 
Dit komt door de ingebouwde capaciteiten om 
TLS-fingerprinting en botdetectie te omzeilen. 

Om binnenkomende data direct te onderscheiden van 
reeds verwerkte concerten zonder valse positieven, 
is een change detection cache geïmplementeerd 
op basis van cryptografische \textit{content hashing}.

Om deze componenten samen te brengen, is een 
modulaire data-architectuur ontworpen die de 
asynchrone verwerking (\texttt{Celery}), de 
database-interactie (\texttt{SQLAlchemy}) en de 
API-logica (\texttt{FastAPI}) strikt scheidt. 

Door extractieregels onder te brengen in databaseconfiguraties in plaats van in de kerncode, 
is de applicatie schaalbaar en eenvoudig uitbreidbaar. 

De geautomatiseerde testsuite bewijst tot slot 
dat het prototype volledig voldoet aan de criteria: 
de matching-engine behaalt een precisie van 100\% 
binnen de datalaag, de verwerking vindt autonoom 
binnen de 4-uurslimiet plaats, en het systeem 
waarborgt een betrouwbare delivery rate via 
automatische retries bij netwerkfalen.

\section*{Kritische reflectie en toekomstig onderzoek}

Doordat de \gls{PoC} technisch is ontworpen en 
functioneel geïmplementeerd via een robuuste 
testsuite, bewijst het de architecturale en 
theoretische haalbaarheid van het concept. 

Voor toekomstig onderzoek liggen er verdere 
mogelijkheden in het uitbreiden van het gebruikersmodel. 
Momenteel focust het ontwerp op expliciete voorkeuren. 
Zoals gesuggereerd door \textcite{Zangerle2018}, kan 
de relevantie verder verhoogd worden door diepere 
contextuele factoren, zoals genre-voorkeuren in 
verhouding tot geografische reisafstand, mee te nemen. 
Dit stelt het systeem in staat om concertsuggesties 
aan te bieden die nóg beter aansluiten bij de 
impliciete behoeften van de concertganger.


